\clearpage
\chapter*{Fe de erratas de la exposición n.\textsuperscript{o} 195}
\addcontentsline{toc}{chapter}{Fe de erratas de la exposición n.\textsuperscript{o} 195}
\setcounter{footnote}{0}
\src{300}

\noindent\textbf{4. TDTE II: El teorema de existencia en la teoría formal de
los módulos.}

\begin{center}
[Seminario Bourbaki, t. 12, 1959/60, n.\textsuperscript{o} 195, 22 p.]
\end{center}

\noindent\textbf{Página 8.} --- La fórmula escrita en la proposición 5.1 solo es
correcta cuando $\Lambda$ es un cuerpo; en el caso general, hay que sustituir
$\mathfrak m_\xi/\mathfrak m_\xi^2$ por el cociente de este espacio por la
imagen de $\mathfrak n_\xi/\mathfrak n_\xi^2$, donde $\mathfrak n$ es el ideal
maximal de $\Lambda$. Además, la definición dada para que $O$ sea liso sobre
$\Lambda$ solo es correcta cuando la extensión residual $k'/k$ es separable.
En el caso general, cf. SGA III, 1.1.

\noindent\textbf{Página 14, observación.} --- Los problemas planteados aquí
quedan completamente resueltos en el caso proyectivo mediante los «esquemas de
Hilbert» (cf. TDTE IV). Ejemplos de Nagata y Hironaka muestran, en cambio, que
los funtores considerados no son necesariamente representables si no se impone
una hipótesis de proyectividad, incluso limitándose a la clasificación de las
subvariedades de dimensión cero de una variedad completa no singular de
dimensión 3; esto está ligado al hecho de que el cuadrado simétrico de una
variedad tal puede no existir.

\src{301}
\noindent\textbf{Página 15, n.\textsuperscript{o} 3.} --- Para un estudio más
completo, véase TDTE V.

\noindent\textbf{Página 16, proposición 3.1.} --- En lugar de «si $f$ es
propio», léase «si $f$ es propio y separable».

\noindent\textbf{Páginas 18 y 19, observaciones, 2.\textsuperscript{o}} --- He
demostrado recientemente que el esquema formal de módulos de una variedad
abeliana sobre un cuerpo es efectivamente liso sobre el anillo de Witt; en
otros términos, todo esquema abeliano sobre un anillo artiniano local que sea
cociente de otro proviene, por reducción, de un esquema abeliano sobre este
último. La demostración utiliza únicamente las propiedades de varianza de la
clase de obstrucción al levantamiento introducida en la exposición
n.\textsuperscript{o} 182, página 12. Recordemos asimismo que los esquemas de
módulos de las curvas de género $g$ y de los esquemas abelianos polarizados han
sido construidos por Mumford (cf. SHMT).

\noindent\textbf{Página 19, párrafo 5.} --- Hay que suponer que la sección
considerada para definir $U$ sea un no divisor de cero no solo sobre
$\mathfrak X$, sino sobre todo $X_n$.

\noindent\textbf{Página 20, líneas 12 y 13.} --- En lugar de
\[
  G_x=G\bigl(\operatorname{Spec}(\widehat{\mathcal O}_{\mathfrak X,x})\bigr),
\]
léase
\[
  G_x=G\bigl(\operatorname{Spec}(\mathcal O_{\mathfrak X,x})\bigr)
  \subset
  G\bigl(\operatorname{Spec}(\widehat{\mathcal O}_{\mathfrak X,x})\bigr),
\]
y, en lugar de «$\mathcal O'_{0x}$», léase
«$\widehat{\mathcal O}'_{0x}$».
